% =========================================================================
%  working_notes.tex -- discussion moved OUT of the compiled supplement.
%
%  WHAT IS HERE.  Remarks that had been promoted from the project's working notes
%  into printed text and that NOTHING in the corpus points at.  They are kept
%  verbatim because several record why a route was DECLINED, which is what a
%  referee asks and a later author re-derives at cost.
%
%  WHY ONLY THREE.  Twenty remarks were identified as working notes in printed
%  form, 38.6 of the supplement's 434 pages.  Seventeen of them carry 174 PRINTED
%  cross-references between them -- rem:pasrc alone is cited 48 times, rem:obligations
%  22, rem:rootnscope 21, rem:lfnonlinear 20 -- so they are load-bearing in the
%  corpus as it stands and cannot be moved without rewriting the sentences that
%  cite them.  That is an authorial decision, not a mechanical one, and it is left
%  open.  These three have ZERO printed references and moved cleanly.
%
%  INPUT BY NO DOCUMENT.  driver.tex inputs eight files and this is not among
%  them, so it contributes to no counter and no label here is live.
%
%  Moved 2026-08-10.
% =========================================================================


%% ---- rem:ffcompare ----
\begin{remark}[The Fève--Florens identification route, and why it does not transplant]\label{rem:ffcompare}
\emph{Amended 2026-08-05.} This remark previously read the transplant of \citet[Prop.~2.2]{feve2014}
to the conditional-independence design as a route that reaches the linearised condition at the cost of
an excess, and priced the excess. \emph{That reading was wrong, and the direction of the error matters}:
the transplanted hypothesis is not stronger than theirs, it is unsatisfiable. Their completeness
condition is posed on $g\mapsto E[g\mid W]$, which does not annihilate constants. Its analogue here
must be posed on $\Lambda\mapsto E[\widetilde W\Lambda\mid\Delta U,\Delta Z]$ ---
\emph{$\IR^{\mathsf L-1}$-valued at every $\mathsf L$, one component per contrast, and scalar at
$\mathsf L=2$; the argument below uses only that the map kills $\Lambda\equiv1$, which is a statement
about every component at once} --- and $\widetilde W$ is
centred given $\Delta Z$, so under A1 that operator kills every $\Lambda$ agreeing almost surely with a
function of $(\Delta U,\Delta Z)$ --- the constants included. The condition fails at $\Lambda\equiv1$
at every law A1 admits, hence at every law A1b(iii) admits. \Cref{lem:nusigma}\textup{(v)} is the
unconditional shadow of the same identity and was in this corpus the whole time.
The corollary that rested on it is withdrawn; see \cref{rem:a1cwithdrawn}. What follows is retained
because the comparison with \citet{feve2014} is still the right way to see \emph{why} the two designs
part company, but it should be read as an account of a closed route, not an available one.

\emph{Notation, fixed once.} \Citet{feve2014} write the endogenous regressors as $Z_1,Z_2$ and the
instrument as $W$; in this paper the endogenous regressors are $X_1,X_2$ and $Z$ is reserved for the
exogenous controls. All quotations below are verbatim \emph{in their notation}; the translation is
$Z_t\mapsto X_t$ throughout, and no symbol inside a quotation has been altered.

\emph{What it delivers.} Their Proposition~2.2 states: ``Let $\varphi\in E$. If $(Z_1,Z_2)$ is strongly
identified by $W$ and if $Z_1$ and $Z_2$ are measurably separated the function $\varphi$ is identified.''
The two primitives would become A1d and A1b(ii) here, and the argument would run the same in outline:
completeness removes the period difference
from inside the conditional expectation, and measurable separability then forces the direction to be
constant, which \cref{def:spaces} sets to zero because the direction lies in $\mathcal V$.
\emph{The outline is where it stops.} The first primitive cannot be assumed in this design, so the
argument has no starting point and A1c gains no primitive sufficient condition from it. Had it been
available it would have converted the draft's condition \texttt{eq-condcomppanel} at \texttt{:127--130}
from an unreconciled display into \cref{prp:kerchar}(c) --- with two corrections: the conditioning set
must include $\Delta Z$, which the draft omits, and the statement is about the linearisation, not about
the structural function.

\emph{What it does not deliver, and why.} It does not reach A1b(iii). In \citet{feve2014} the
restriction is a conditional mean and the operator is linear in $\varphi$, so injectivity of the operator
\emph{is} identification. Under conditional independence the operator equation is nonlinear in $\varphi$,
and the two statements separate. The obstruction is nameable: membership of $\mathcal V_{CI}$ is the
family $E[\widetilde W\varsigma(\Delta Y-\Delta\varphi)\mid\Delta Z]=0$ indexed by test functions of the
\emph{observed} residual at the trial value, whereas running the argument of \cref{prp:kerchar} globally
would require the strictly larger family
$E[\widetilde W\varsigma(\Delta Y-\Delta\varphi)\varkappa(\Delta U)\mid\Delta Z]=0$, indexed also by test
functions of the \emph{true} residual, about which the restriction says nothing at a trial $\varphi$. The
two families coincide at $\varphi=\vdagger$ and nowhere else, which is exactly the gap between
first-order and global identification. \Cref{prp:compact} closes off the obvious repair: because
$\mathcal A_{\mathrm{cond}}$ is compact it is not bounded below, so A1c does not even deliver a
neighbourhood of uniqueness.

\emph{Why recording the gap is the right response, and the precedent is sharper than a difficulty claim.}
\citet{feve2014} reject the hypothesis of their own Proposition~2.2, and their stated reason is the
excess A1d reproduces exactly: ``The assumption of Proposition 2.2 is too strong because it assumes that
$E(g(Z_1,Z_2)\mid W)=0$ implies $g(Z_1,Z_2)=0$ a.s.\ for any $g$, and we only use this property for
$g(Z_1,Z_2)$ of the form $\varphi(Z_2)-\varphi(Z_1)$.'' A1d has the same shape: it quantifies over every
$\Lambda\in L^{2}(P_{X_1,X_2})$, while the corollary consumed it only at $\Lambda=\Delta h$, a
difference. \emph{Both papers reach a completeness condition stated over a class strictly larger than the
one the argument uses, and both decline to repair it by restricting the class} --- because injectivity of
$\Lambda\mapsto E[\widetilde W\Lambda\mid\Delta U,\Delta Z]$ on differences alone is no longer a
recognisable completeness condition and has no primitive characterisation of its own.
\emph{The shapes coincide and the consequences do not.} Their excess is an excess: their condition is
satisfiable and merely asks more than the argument consumes. Ours was not, and restricting the class
would have been the only way to state anything satisfiable at all --- which is why the condition is
withdrawn rather than weakened.

\emph{The two asymmetries are not the same asymmetry, and the synthesis has to say so.} Theirs runs from
the assumption to its use: the hypothesis is \emph{stronger than the argument consumes}. The one recorded
in the previous paragraph runs the other way: the restriction supplies \emph{less information than the
argument needs}, the $\varsigma$ family against the $(\varsigma,\varkappa)$ family. They are opposite
inequalities and they meet at one point --- \emph{the recognisable completeness conditions available in
either design are stated over a coarser class than the design's own information}, so an assumption
calibrated to the class is too strong and the information calibrated to the class is too weak. \emph{That
is where a future improvement would have to act, in either paper}: not by strengthening the completeness
condition, but by finding a restriction on $P$ stated over exactly the class the argument uses.
\emph{In this design that is not an improvement but the only possibility}, since the completeness
condition stated over the coarser class is not merely too strong here --- it is unsatisfiable, the
operator having the constants in its kernel. A1b(iii) is accordingly kept as an assumption with no
primitive route offered, and the transplanted condition is recorded as closed rather than as taken.
\end{remark}


%% ---- rem:draftrec ----
\begin{remark}[Two features of the linearisation that the derivation must not absorb]\label{rem:draftrec}
\emph{(i) The operator's arguments.} In \eqref{eq:Mdef} the $(x_1,x_2)$-integration is an expectation,
the $W$-argument is replaced by the contrast index $l$ --- legitimate because the $\mathsf L-1$
contrasts \eqref{eq:contrast} span the conditionally centred functions of $W$, so nothing in the $W$
direction is approximated --- and
the $z$-argument is removed by the Bierens index $t$. What is left is indexed by $(u,t,l)$ with $u$
scalar and $l$ ranging over a \emph{finite} set, which is why the existing Landweber--Fridman machinery
carries over with an enlarged index set
and no new curse of dimensionality. \emph{At $\mathsf L=2$ the index takes one value and the object is
indexed by $(u,t)$ as earlier rounds printed.}

\emph{(ii) Conditional versus unconditional centring.} A raw contrast between the $W$-cells,
unconditional in $Z$, is \emph{not} the conditional operator reweighted by a scalar, and the two do not
share a range or a kernel. \emph{The correct comparison is
between kernels.} Write
$f_j(u):=f_{\Delta U\mid X_1,X_2,\Delta Z,W}(u\mid\cdot,\cdot,\cdot,j)$, the same conditional density
that appears in \eqref{eq:kernel}, and $p_{X,j}:=P(W=j\mid X_1,X_2,\Delta Z)$,
$j\in\{0,\dots,\mathsf L-1\}$. The kernel of
$\mathcal A_{\mathrm{cond}}$ in \eqref{eq:kernel} is built from
\begin{equation}\label{eq:kernelcompare}
E\big[\widetilde W_l\,f_{\Delta U\mid X_1,X_2,\Delta Z,W}(u)\,\big|\,X_1,X_2,\Delta Z\big]
\;=\;p_{X,l}\,f_l(u)\;-\;\varrho_l\sum_{j=0}^{\mathsf L-1}p_{X,j}\,f_j(u),
\end{equation}
whereas the unconditional operator is built from the raw cell differences $f_l(u)-f_0(u)$ up to sign.
\emph{The comparison is not the binary one, and the $\mathsf L\ge3$ correction is a matter of
\textbf{which} differences appear and not only of a scalar factor.} Suppose $p_{X,j}=\varrho_j$ for
every $j$ --- equivalently $W\upmodels(X_1,X_2)\mid\Delta Z$, which says that \emph{the instrument does
not shift the endogenous regressor given $\Delta Z$}. Then, since $\sum_j\varrho_j=1$, the right-hand
side of \eqref{eq:kernelcompare} becomes
\begin{equation}\label{eq:kernelcompareci}
\varrho_l\Big\{f_l(u)-\sum_{j=0}^{\mathsf L-1}\varrho_jf_j(u)\Big\}
\;=\;\varrho_l\sum_{j=0}^{\mathsf L-1}\varrho_j\big\{f_l(u)-f_j(u)\big\} ,
\end{equation}
a $\varrho$-weighted combination of \emph{all} $\mathsf L-1$ raw differences available at $l$ --- cell
$l$ against the $\varrho$-mixture of every cell --- and \textbf{not} a multiple of any single one.
At $\mathsf L=2$ the sum has one nonvanishing term and \eqref{eq:kernelcompareci} is
$\varrho_1\varrho_0\{f_1(u)-f_0(u)\}$, so there, and only there, do the two operators coincide up to the
scalar $\varrho_1\varrho_0$, and there conditional independence is also \emph{necessary} for that
coincidence, $p_{X,1}\varrho_0f_1-\varrho_1p_{X,0}f_0$ being proportional to $f_1-f_0$ exactly when
$p_{X,1}=\varrho_1$. \emph{At $\mathsf L\ge3$ conditional independence does not deliver it}: what
survives $p_X=\varrho$ is the choice of comparison cell, and \eqref{eq:kernelcompareci} is a scalar
multiple of $f_l-f_0$ only under a further coincidence among the cell densities --- $f_j=f_0$ for every
$j\notin\{0,l\}$ suffices and returns $\varrho_l(1-\varrho_l)\{f_l(u)-f_0(u)\}$ --- which is a
restriction on the model and not on the instrument. \emph{Conditional independence must not be assumed in
any case, and no equality of ranges or of kernels is claimed here.} The identity
$E[\widetilde W_l\,g(W)\mid\Delta Z]=\varrho_l\{g_l-\sum_{j}\varrho_jg_j\}$, which is correct, holds at
$\sigma(\Delta Z)$ and does not extend to $\sigma(X_1,X_2,\Delta Z)$; see convention (P4b) and its scope
warning in \texttt{proofs.tex} --- \emph{(P4b) states it there in \textbf{this same general form}, at
every $\mathsf L$ and cell by cell: (P4b)'s display is the mixture one, and the two-cell
$\varrho(1-\varrho)\{g_1-g_0\}$ appears beside it only inline, as the $\mathsf L=2$ reduction. What
separates (P4b) from \eqref{eq:kernelcompareci} is therefore \textbf{not} the number of cells --- the two
right-hand sides are one object --- but the conditioning: (P4b) is the identity at $\sigma(\Delta Z)$,
where the cell weights are the $\varrho_j$, whereas \eqref{eq:kernelcompare} is the corresponding object
at $\sigma(X_1,X_2,\Delta Z)$, where they are the $p_{X,j}$ and where (P4b) is unavailable, and
\eqref{eq:kernelcompareci} is what the latter collapses to when the two sets of weights agree. That is
the whole content of this clause}.
% RETYPED, PHASE 5a ROUND 2, under a named grant into this otherwise closed file.
% This clause described (P4b) as stating the identity "in the two-cell form
% varrho(1-varrho){g_1 - g_0}" and referred to "(P4b)'s two-cell display".  Both
% were true before phase 5a and false after it: (P4b) now states the general
% mixture form at every L and carries no two-cell display.  The clause was thus a
% claim about proofs.tex that proofs.tex contradicted -- and proofs.tex, at
% lem:pbinary, was simultaneously printing that this very deferral had been
% DISCHARGED.  Both ends are corrected in one pass.  The sigma-field distinction,
% which was and remains the load-bearing content, is preserved verbatim in
% substance.
The weights $p_{X,j}$ and $\varrho_l$ in \eqref{eq:kernelcompare} are
bounded for free, every factor lying in $[0,1]$, and \emph{no prefactor is inverted anywhere in this
development}: the one argument that used to divide by $\varrho(1-\varrho)$ --- \cref{lem:quantu}'s
extension from a countable set of $u$ to all of $\IR$ --- now propagates a zero along a right-continuous
function instead, bounding nothing and inverting nothing, so A4(i) is \textbf{not} among that lemma's
hypotheses and no case split on the degenerate region is made there.
\emph{That \cref{lem:quantu} has no proof-step consumer, and why it is kept notwithstanding, are recorded
in its own statement} --- the single home for that status, on \cref{lem:nusigma}\textup{(ii)}'s model.
% RETYPED, PHASE 5a ROUND 2, under the same named grant.  This sentence read: "what
% A4(i) delivers is invertibility of the prefactor in the one place the prefactor
% actually appears, namely at sigma(Delta Z) in lem:quantu."  Phase 5a round 1
% repaired lem:quantu's Step 3 and the repaired proof prints the OPPOSITE -- "no
% division by a prefactor" -- and A4(i) is not a hypothesis of lem:quantu at all.
% ROUND-2's NOTE ON PLACEMENT IS SPENT, AND PHASE 5c's SUCCESSOR NOTE IS NOW SPENT
% TOO.  Round 2 put the consumer-status note HERE only because its grant did not
% reach the lemma body.  Phase 5c owned the body and put the note there, on
% lem:nusigma(ii)'s model, but RETAINED the paragraph here as well and reported the
% residual to the critic as a duplication rather than a contradiction.
% PHASE 4c, 2026-08-10, CUTS IT, and the judgement is recorded because the
% alternative was defensible.  THE GROUND IS NOT TIDINESS.  (1) Nothing in the cut
% paragraph was unique: all three of its facts -- no consumer, why kept, not
% available as an inherited step -- are in the lemma body, which is STRICTLY more
% informative, naming rem:localise(1) as the replacement route where this paragraph
% said only "a cell-factorisation route".  (2) The duplicated content is a DATED
% STATUS FACT ("as of 2026-08-09").  A dated status printed twice, ~320 lines apart,
% is a stale-claim generator: the pass that eventually gives lem:quantu a consumer
% will update one copy and leave the other, and no string census finds that -- it is
% the corpus's dominant defect species, manufactured on a timer.  (3) The reader is
% not stranded and the lemma body's own claim stays TRUE: eq:kernelcompareci's
% clause still cites lem:quantu in the sentence immediately above, so it still
% "sends a reader there", and the pointer left behind says where the status lives.
% BINARY GATE 1 IS UNAFFECTED: lem:quantu's statement and its proof both exist and
% are matched by label.  "No consumer" is a fact about which STEPS cite it, not
% about whether it is proved.
There is in any case no second operator to reconcile with:
$\mathcal A_{\mathrm{cond}}$ \emph{is} the operator governing $\hat\varphi$.

\emph{(iii) The adjoint, and the factor that must not be cancelled.} The adjoint itself is
\begin{equation}\label{eq:adjointdisplay}
\begin{split}
\big(\mathcal A^\ast_{\mathrm{cond}}\ell\big)(x)\;&=\;\sum_{l=1}^{\mathsf L-1}\varpi_{\mathsf w}(l)
\int_{\IR\times\mathcal T}\mathsf k(u,t,l;x)\,\ell(u,t,l)\,\Pi_u(du)\,\Pi_t(dt)\\
&=\;\frac{\nu_{\widetilde W'\,(\mathcal W\ell)}(x)}{\pi(x)}
\qquad\text{with }\mathsf k\text{ from \eqref{eq:kernel}},
\end{split}
\end{equation}
so that $\langle\mathcal A_{\mathrm{cond}}h,\ell\rangle_{L^2(\Pi)}=\langle h,\mathcal A^\ast_{\mathrm{cond}}\ell\rangle_{L^2(\pi)}$
holds by construction; this is \cref{lem:adjoint}. \emph{Two things about
\eqref{eq:adjointdisplay} are load-bearing.} First, in the \emph{pairing}
$\langle\mathcal A h,\ell\rangle$ the $L^2(\pi)$ inner product contributes a factor $\pi$ and the
adjoint's own formula contributes $\pi^{-1}$, and the two cancel; that is a trap, because a reader who
cancels them concludes that the $L^2(\pi)$ weighting is immaterial, when the surviving $\pi^{-1}$ in
\eqref{eq:adjointdisplay} is exactly why the source condition A2 is imposed on $\nu/\pi$ and not on
$\nu$. Second, both subscripts $(\theta,\varphi)$ of the density are substantive:
\cref{prp:frechet}(ii) shows the $\varphi$-dependence is real, the derivative map being Lipschitz in
$\varphi$ with a nonzero constant; and \cref{rem:R1}(iv) records that the operator's singular values,
degree of ill-posedness and source index are all $\theta$-specific, because \eqref{eq:Adef} evaluates
the conditional density at a $\theta$-dependent point. \emph{What is $\theta$-free is a different object
--- the solution set of \eqref{eq:opeq} in $\varphi$, by \cref{lem:shiftinv}. The operator is not
$\theta$-free; only its zero set is.}

\emph{(iv) Two symbols.} The kernel $\un(x\ge v)$ is the kernel of
$\mathcal D^{-1}$, the integration operator of the derivative parameterisation; with the analysis
conducted in levels it does not appear. The letter $\psi$ is reserved for the moment function, so a
generic element of the range space is written $\ell$. \emph{And the range measure may not be left
implicit}: without $\Pi$ there is nothing to fix the range space, and displays that differ in whether a
$W$-argument survives silently use different ones. Under \cref{def:operator} the contrast vector
$\widetilde W$ replaces the $W$-argument by the contrast index $l$ once and for all, $\Pi$ gains
$\Pi_{\mathsf w}$ as its third factor, and the range is $L^{2}(\Pi)$ throughout.
\end{remark}


%% ---- rem:nophifree ----
\begin{remark}[What the impossibility result says, and where it has been seen before]\label{rem:nophifree}
\Cref{prp:nophifree} is the paper's negative result, and three readings of \eqref{eq:tilted} are worth
separating.

\emph{(i) The condition, in words.} Write $\tilde f^{\,g}_j(x):=E[g\mid X_j=x]f_{X_j}(x)$ for the
$g$-tilted density of the endogenous regressor in period $j$ --- a signed density of total mass $E[g]$.
Condition \eqref{eq:tilted} says $\tilde f^{\,g}_1=\tilde f^{\,g}_2$: \emph{the two periods' $g$-tilted
marginals coincide.} At $g\equiv1$ this is marginal stationarity of $X$, so even the unweighted moment
$E[\Delta U]$ is free of $\varphi$ only when the regressor's marginal law does not move between the two
periods. At $g=\Delta Z_j$ it is $\nu\equiv0$, which is probe CP-1 and the design in which the debiasing
weight degenerates to zero and the moment reduces to the differenced least-squares one. The failure of
\eqref{eq:tilted} is therefore not a pathology to be assumed away: it is what makes the panel design
informative in the first place, and a design in which it holds for the weights of interest is one in
which there is nothing to debias.

\emph{(ii) The same condition, reached independently in the linear case.} In the first-differenced panel
with a conditional-mean restriction, \citet{feve2014} observe that the operator
$\varphi\mapsto E[\varphi(X_2)-\varphi(X_1)\mid W]$ annihilates \emph{every} $\varphi$ when
$(X_1,X_2,W)$ has the same law as $(X_2,X_1,W)$ --- exchangeability of the two periods given the
instrument, which in particular equates the laws of $(X_1,W)$ and $(X_2,W)$. That is \eqref{eq:tilted}
imposed simultaneously for every bounded $g=g(W)$, since
$\nu_{g(W)}(x)=\sum_{j=0}^{\mathsf L-1}g(j)\{f_{X_2,W}(x,j)-f_{X_1,W}(x,j)\}$. Their Proposition~2.1 gives the
adjoint of that operator as an integral against the same difference of joint densities over the same
reference density; see \cref{rem:adjointreconcile}(iv). \emph{The difference structure is therefore not a
device of this paper: it is inherited from first differencing and is shared with the linear
conditional-mean design.} What \cref{prp:nophifree} adds is the one-weight-at-a-time equivalence --- their
condition is the simultaneous one, and is correspondingly stronger --- and the enlargement of $g$ from
functions of $W$ to functions of $(W,\Delta Z)$, which is the class the debiasing construction actually
searches.

\emph{(iii) What is new here, stated precisely.} The sharp difference from \citet{feve2014} is not that
their operator is linear and ours is not, which is true but too coarse. It is that
$\mathcal A_{\mathrm{cond}}$ is a \emph{Fr\'echet derivative of a nonlinear map} rather than a fixed
linear conditional-expectation operator, and three consequences follow that have no counterpart there:
the operator depends on the point of linearisation, with the Lipschitz modulus of
\cref{prp:frechet}(ii); the sample operator is noisy \emph{as an operator}, which is why the
regularisation conditions of this block --- A7, and the retired A6 before it --- are stated for
$\hat T$ and not for $T$; and the moment carries a second-order term in $\varphi$, which is
\cref{cor:udep}'s subject and which a linear operator equation does not produce.

\emph{(iv) Why the negative result creates the rate requirement, and what the chain does not contain.}
\Cref{cor:udep} closes the chain. If the debiasing weight ignored the residual, the moment would reduce
to a $\sigma(W,\Delta Z)$-measurable weight times the residual, and orthogonality would then force
\eqref{eq:tilted}; in any design where \eqref{eq:tilted} fails for such weights, the selected
$\ell_\star$ must vary with the residual. That is what puts $\partial_u\ell_\star$ into the second-order
term, and it is why a quadratic remainder in $\varphi$ is unavoidable within this moment class rather
than an artefact of the expansion. \emph{Two things are not in the chain and should not be read into it.}
First, $\partial_u\ell_\star\not\equiv0$ is \emph{necessary} for the second-order form to be nonzero and
not sufficient: the conditional expectation
$q(x_1,x_2)=E[\partial_u\ell_\star\,\widetilde W\mid X_1=x_1,X_2=x_2]\in\IR^{p}$ could still vanish by cancellation,
and excluding that is a separate hypothesis. Second, \cref{prp:nophifree} is an equivalence, not an
impossibility theorem: it does not assert that \eqref{eq:tilted} fails. Whether it fails is a property of
the application, testable through the Picard diagnostic and through $\hat\nu_g$ directly, and the paper
should report it rather than assert it.

% EXTRACTED FROM THE STATEMENTS OF prp:nophifree AND cor:a1cprim on 2026-08-02.
% Both carried a "not a hypothesis" paragraph inside the printed statement; both
% paragraphs are about the hypothesis LIST and not about what is claimed.  The
% clause is struck from neither list -- it was never in either -- so nothing about
% either result changes.
\emph{(v) A1b(i) is \textbf{not} a hypothesis of \cref{prp:nophifree}, and it is not one of
the withdrawn corollary either.} For \cref{prp:nophifree}: the proof's third step upgrades
``$\nu_g/\pi\perp\mathcal V$'' to ``$\nu_g\equiv0$'' using
$L^{2}(\pi)=\mathcal V\oplus\{\text{constants}\}$ --- which is \cref{def:spaces} --- together with
$\int_{\mathcal X}\nu_g\,dx=0$ from \cref{lem:nusigma}(iv); $\vdagger$ occurs nowhere in the statement or
the proof, and the map is quantified over generic $\varphi,\varphi_0\in\mathcal V$. For
the withdrawn corollary of \cref{rem:a1cwithdrawn}: the ``constant, hence zero'' step of \cref{lem:nusigma}\textup{(ii)} is
\cref{def:spaces} and not that clause, and the conclusion is about a generic
$h\in\ker\mathcal A_{\mathrm{cond}}\subseteq\mathcal V$.
% ARCHAEOLOGY.  A1b(i) was listed as a hypothesis of both through round 7, and is
% struck from both; cor:a1cprim inherited it from lem:nusigma(ii).
\end{remark}
